\documentclass[11pt]{article}

% ============================================================
% Packages
% ============================================================
\usepackage[margin=1in]{geometry}
\usepackage{amsmath,amssymb,mathtools,amsthm}
\usepackage{enumitem}
\usepackage{hyperref}
\usepackage{xcolor}
\usepackage{microtype}
\usepackage{tikz}

% ============================================================
% Macros
% ============================================================
\newcommand{\R}{\mathbb{R}}
\newcommand{\Pbb}{\mathbb{P}}
\newcommand{\Qbb}{\mathbb{Q}}
\newcommand{\E}{\mathbb{E}}
\newcommand{\dd}{\,\mathrm{d}}
\newcommand{\F}{\mathcal{F}}

% ============================================================
% Environments
% ============================================================
\theoremstyle{plain}
\newtheorem{theorem}{Theorem}[section]
\newtheorem{proposition}[theorem]{Proposition}

\theoremstyle{definition}
\newtheorem{definition}[theorem]{Definition}
\newtheorem{example}[theorem]{Example}
\newtheorem{exercise}[theorem]{Exercise}

\theoremstyle{remark}
\newtheorem{remark}[theorem]{Remark}

% ============================================================
% Title
% ============================================================
\title{Option Pricing --- Lecture Notes}
\author{Nathan Sauldubois}
\date{MSFE6233 --- Refresher}

\begin{document}

\maketitle
\vspace{-1em}

\noindent\textbf{Lecture title.}
\emph{Brownian Motion, Stochastic Calculus, Girsanov Theorem, Option Pricing and Delta Hedging}

\vspace{0.5em}
\hrule
\vspace{1em}

\noindent\textbf{Goals of the lecture.}
\begin{itemize}[leftmargin=1.4em]
    \item A few refreshers on key notions we have seen throughout the lectures.
\end{itemize}

\tableofcontents
\bigskip

% ============================================================
\section{Brownian Motion}
% ============================================================


\begin{exercise}[Compute $\int_0^t e^{B_s}\,\mathrm{d} B_s$ without $\mathrm{d} B$]
\begin{enumerate}
\item Apply It\^o's formula to $e^{B_t}$ and isolate the term $\int_0^t e^{B_s}\,\mathrm{d} B_s$.
\item Deduce an expression for $\int_0^t e^{B_s}\,\mathrm{d} B_s$ involving only $e^{B_t}$ and ordinary integrals.
\end{enumerate}
\end{exercise}

\begin{exercise}[Compute $e^{\int_0^t f(s)\,\mathrm{d} B_s}$ without $\mathrm{d} B$]
\begin{enumerate}
\item Apply It\^o's formula to the exponential martingale associated with $\int_0^t f(s)\,\mathrm{d}B_s$.
\item Deduce an expression involving only ordinary integrals.
\end{enumerate}
\end{exercise}

\section{Girsanov}

\begin{exercise}[Girsanov and the inverse Gaussian law]
Let $(B_t)_{t\ge 0}$ be a standard Brownian motion under a probability measure $\mathbb{P}$.
Fix $\mu>0$, $\sigma>0$, and $a>0$, and define the drifted Brownian motion
\[
X_t := \mu t + \sigma B_t.
\]
Let
\[
\tau_a := \inf\{t\ge 0 : X_t = a\}
\]
be the first hitting time of level $a$.

\begin{enumerate}[leftmargin=1.5em, label=(\alph*)]

\item Show that $\tau_a<\infty$ almost surely.

\item Define
\[
\theta := \frac{\mu}{\sigma},
\qquad
Z_t := \exp\left(-\theta B_t - \frac12 \theta^2 t\right).
\]
Show that $(Z_t)_{t\ge 0}$ is a martingale, and define a new probability measure $\mathbb{Q}$ on $\mathcal{F}_t$ by
\[
\frac{\mathrm{d}\mathbb{Q}}{\mathrm{d}\mathbb{P}}\Big|_{\mathcal{F}_t}
= Z_t.
\]
Use Girsanov's theorem to show that
\[
\widetilde B_t := B_t + \theta t
\]
is a Brownian motion under $\mathbb{Q}$.

\item Rewrite $X_t$ under $\mathbb{Q}$ in terms of $\widetilde B_t$, and deduce that
\[
X_t = \sigma \widetilde B_t.
\]
Thus, under $\mathbb{Q}$, $\tau_a$ is the first hitting time of level $a/\sigma$ by a standard Brownian motion.

\item Recall, or prove using the reflection principle, that if
\[
\tau_b^0 := \inf\{t\ge 0 : \widetilde B_t=b\},
\qquad b>0,
\]
then $\tau_b^0$ has density
\[
f_{\tau_b^0}(t)
=
\frac{b}{\sqrt{2\pi t^3}}
\exp\left(-\frac{b^2}{2t}\right),
\qquad t>0.
\]

\item Use the change of measure to compute the density of $\tau_a$ under $\mathbb{P}$.

\textit{Hint:} for events in $\mathcal{F}_t$,
\[
\mathbb{P}(A)=\mathbb{E}^{\mathbb{Q}}\left[Z_t^{-1}\mathbf 1_A\right].
\]
On the event $\{\tau_a=t\}$, use the identity
\[
B_t = \frac{a-\mu t}{\sigma}.
\]

\item Deduce that $\tau_a$ has density
\[
f_{\tau_a}(t)
=
\frac{a}{\sigma\sqrt{2\pi t^3}}
\exp\left(
-\frac{(a-\mu t)^2}{2\sigma^2 t}
\right),
\qquad t>0.
\]
This is the density of an inverse Gaussian distribution.

\end{enumerate}
\end{exercise}


\section{Martingale Pricing}

% ------------------------------------------------------------
\begin{exercise}[Asian option (arithmetic average): martingale pricing, state augmentation, and PDE.]
Assume constant $(r,\sigma)$ and under $\mathbb{Q}$:
\[
\frac{\mathrm{d}S_t}{S_t}=r\,\mathrm{d}t+\sigma\,\mathrm{d}B_t,
\qquad S_0>0.
\]
Consider the arithmetic-average Asian call with payoff
\[
H=\left(\frac1T\int_0^T S_u\,\mathrm{d}u-K\right)^+.
\]
Define the running integral
\[
I_t:=\int_0^t S_u\,\mathrm{d}u.
\]

\begin{enumerate}[leftmargin=1.2em, label=(\alph*), itemsep=0.3em]
\item Write the risk-neutral pricing formula for the option price
\[
V_t=\E^\mathbb{Q}\!\left[e^{-r(T-t)}\left(\frac1T\int_0^T S_u\,\mathrm{d}u-K\right)^+\Bigm|\F_t\right].
\]

\item Show that
\[
\int_0^T S_u\,\mathrm{d}u=I_t+\int_t^T S_u\,\mathrm{d}u,
\]
and deduce that the price can be written as a deterministic function of \emph{two} state variables:
\[
V_t=v(t,S_t,I_t).
\]

\item Derive the dynamics of $(S_t,I_t)$ under $\mathbb{Q}$:
\[
\mathrm{d}S_t=rS_t\,\mathrm{d}t+\sigma S_t\,\mathrm{d}B_t,
\qquad
\mathrm{d}I_t=S_t\,\mathrm{d}t.
\]
Explain why $(S_t,I_t)$ is Markov even though $S$ alone is not enough for pricing this payoff.

\item Assume $v\in C^{1,2,1}$ is smooth enough. Apply It\^o's formula to $v(t,S_t,I_t)$ and derive the PDE satisfied by $v$:
\[
\partial_t v + rs\,\partial_s v + s\,\partial_i v + \frac12\sigma^2 s^2\,\partial_{ss}v - rv=0,
\]
with terminal condition
\[
v(T,s,i)=\left(\frac{i}{T}-K\right)^+.
\]

\item Define the discounted price $\bar V_t=e^{-rt}V_t$. Show directly from the risk-neutral pricing formula that $(\bar V_t)$ is a $\mathbb{Q}$-martingale.

\item (Numerical / implementation-oriented) Propose a Monte Carlo estimator for $V_0$ based on an Euler discretization of
\[
\int_0^T S_u\,\mathrm{d}u
\approx \sum_{k=0}^{n-1} S_{t_k}\Delta t.
\]
State clearly:
\begin{itemize}[leftmargin=1.2em]
\item what is simulated under $\mathbb{Q}$,
\item what is averaged,
\item where the discount factor appears.
\end{itemize}

\item (Bonus, medium-hard) Let
\[
G:=\exp\!\left(\frac1T\int_0^T \log S_u\,\mathrm{d}u\right)
\]
be the geometric average. Explain why the geometric-average Asian option is analytically tractable in Black--Scholes, and suggest using it as a \emph{control variate} for Monte Carlo pricing of the arithmetic Asian option.
\end{enumerate}
\end{exercise}

\begin{exercise}[Continuous-time lookback digital (one-touch) in Black--Scholes]
We work under the risk-neutral measure $\mathbb{Q}$ using the money-market numeraire.
Assume the stock pays a continuous dividend yield $q\ge 0$ and follows
\[
\frac{\dd S_t}{S_t}=(r-q)\,\dd t+\sigma\,\dd B_t,
\qquad S_0>0,\ \sigma>0.
\]
Fix a barrier level $H>S_0$ and a maturity $T>0$.

Consider the \emph{lookback digital (one-touch)} payoff
\[
\Pi := \mathbf 1_{\left\{\sup_{0\le t\le T} S_t \ge H\right\}}.
\]
This contract pays \$1 at time $T$ if the stock \emph{ever} touches or exceeds $H$ before $T$.

\medskip
\noindent\textbf{Goal.} Compute the no-arbitrage price
\[
V_0 = \E^{\mathbb{Q}}\!\left[e^{-rT}\,\Pi\right]
= e^{-rT}\,\mathbb{Q}\!\left(\sup_{0\le t\le T} S_t \ge H\right).
\]

\begin{enumerate}[leftmargin=1.6em, itemsep=0.35em]

\item \textbf{Log-transform and reduction to a hitting event.}
Define the log-return process
\[
X_t := \ln\!\Big(\frac{S_t}{S_0}\Big).
\]
\begin{enumerate}[leftmargin=1.6em, itemsep=0.25em]
\item Show by It\^o's formula that
\[
X_t = \Big((r-q)-\tfrac12\sigma^2\Big)t + \sigma B_t
=: \mu t + \sigma B_t,
\qquad \mu:=(r-q)-\tfrac12\sigma^2.
\]
\item Show that the barrier event can be rewritten as
\[
\left\{\sup_{0\le t\le T} S_t \ge H\right\}
=
\left\{\sup_{0\le t\le T} X_t \ge a\right\},
\qquad
a:=\ln\!\Big(\frac{H}{S_0}\Big)>0.
\]
\end{enumerate}

\item \textbf{A key lemma: hitting probability for a drifted Brownian motion.}
Let $(B_t)_{t\ge0}$ be a standard Brownian motion and define
\[
Y_t := \mu t + \sigma B_t,
\qquad \mu\in\R,\ \sigma>0.
\]
Prove the identity
\begin{equation}\label{eq:hitting-drifted}
\mathbb{Q}\!\left(\sup_{0\le t\le T} Y_t \ge a\right)
=
\Phi\!\left(\frac{\mu T-a}{\sigma\sqrt{T}}\right)
+
\exp\!\left(\frac{2\mu a}{\sigma^2}\right)
\Phi\!\left(\frac{-\mu T-a}{\sigma\sqrt{T}}\right),
\qquad a>0,
\end{equation}
where $\Phi$ is the standard normal CDF.

\medskip
\noindent\textbf{Guided proof.}
\begin{enumerate}[leftmargin=1.6em, itemsep=0.25em]
\item First treat the driftless case $\mu=0$:
show that
\[
\mathbb{Q}\!\left(\sup_{0\le t\le T} \sigma B_t\ge a\right)
=
2\,\mathbb{Q}(\sigma B_T \ge a)
=
2\left[1-\Phi\!\left(\frac{a}{\sigma\sqrt{T}}\right)\right].
\]
Hint: use the reflection principle for Brownian motion.

\item Now handle $\mu\neq 0$ using a \emph{change of measure}.
Define $\theta:=\mu/\sigma$ and the density
\[
Z_T:=\exp\!\left(-\theta B_T - \tfrac12\theta^2 T\right).
\]
Under the tilted measure $\widetilde{\mathbb{Q}}$ defined by
$\dd\widetilde{\mathbb{Q}}=Z_T\,\dd\mathbb{Q}$ on $\F_T$,
the process $\widetilde B_t:=B_t+\theta t$ is a Brownian motion.
Rewrite $Y_t$ under $\widetilde{\mathbb{Q}}$ in terms of $\widetilde B$ and reduce the problem to
a driftless maximum event under $\widetilde{\mathbb{Q}}$.

\item Carefully transform back to $\mathbb{Q}$ using
$\E^{\mathbb{Q}}[\cdot]=\E^{\widetilde{\mathbb{Q}}}[Z_T^{-1}\cdot]$
and simplify to obtain \eqref{eq:hitting-drifted}.
\end{enumerate}

\item \textbf{Closed-form price of the one-touch lookback digital.}
Use \eqref{eq:hitting-drifted} with $Y=X$, i.e.\
$\mu=(r-q)-\tfrac12\sigma^2$ and $a=\ln(H/S_0)$, to deduce
\[
V_0
=
e^{-rT}\left[
\Phi\!\left(\frac{\mu T-a}{\sigma\sqrt{T}}\right)
+
\exp\!\left(\frac{2\mu a}{\sigma^2}\right)
\Phi\!\left(\frac{-\mu T-a}{\sigma\sqrt{T}}\right)
\right],
\qquad
\mu=(r-q)-\tfrac12\sigma^2,\quad a=\ln\!\Big(\frac{H}{S_0}\Big).
\]

\textbf{Sanity checks.}
\begin{itemize}[leftmargin=1.6em, itemsep=0.2em]
\item If $H\downarrow S_0$, so $a\downarrow0$, then $V_0\to e^{-rT}$.
\item If $H\uparrow\infty$, so $a\to\infty$, then $V_0\to 0$.
\end{itemize}

\item \textbf{Optional Monte Carlo verification with discrete monitoring.}
Simulate $S$ on a grid $0=t_0<\cdots<t_N=T$ and approximate
\[
\mathbf 1_{\{\sup_{0\le t\le T} S_t\ge H\}}
\approx
\mathbf 1_{\{\max_{0\le k\le N} S_{t_k}\ge H\}}.
\]
Compute the Monte Carlo estimator
\[
\widehat V_0^{(N)} := e^{-rT}\,\frac1M\sum_{m=1}^M
\mathbf 1_{\{\max_k S_{t_k}^{(m)}\ge H\}},
\]
and compare it to the closed-form value from Question 3.
Comment on the bias as $N$ increases.
\end{enumerate}
\end{exercise}

\section{American Options}

% ------------------------------------------------------------
\begin{exercise}[American vs European put in a three-step tree]
Consider a three-period binomial model with
\[
S_0=100,\qquad u=1.1,\qquad d=0.9,\qquad 1+r=1.02,
\]
and strike $K=100$.

\begin{enumerate}[leftmargin=1.6em]
\item Compute the price of the corresponding European put.

\item Compute the price of the corresponding American put.

\item Compare the two prices and compute the early exercise premium.

\item Identify the nodes where immediate exercise is optimal for the American put.
\end{enumerate}
\end{exercise}




\section{Numerics}

Add a recap of the Euler scheme.

\end{document}